\section{Conclusão}
\label{sec:conclusao}

Este trabalho especificou, modelou e implementou um validador de
placas de veículos brasileiras. O problema foi tratado como o
reconhecimento de uma linguagem regular, resolvido por autômatos
finitos determinísticos.

\subsection{O que foi entregue}

\begin{itemize}
  \item uma especificação objetiva do que é uma placa válida, com
        as decisões de projeto registradas (Seção~\ref{sec:especificacao});
  \item a definição formal de três autômatos, um para o padrão
        anterior ao Mercosul, um para o Mercosul e um para o
        Mercosul de placas convertidas (Seção~\ref{sec:modelagem});
  \item os diagramas de transição (Seção~\ref{sec:diagramas});
  \item a demonstração de que cada autômato é mínimo, e a redução
        do validador combinado de 14 para 11 estados (Seção~\ref{sec:minimizacao});
  \item uma implementação em Python, verificada com casos válidos e
        inválidos e comparada com a expressão regular em
        $29\,524$ cadeias, sem divergência (Seção~\ref{sec:testes}).
\end{itemize}

\subsection{O que se aprendeu}

A modelagem mostrou que o problema tem uma estrutura simples: a
placa é uma sequência de etapas, e cada estado do autômato
representa uma delas. Isso explica por que os autômatos dos dois
padrões são mínimos por si sós. A única redução aconteceu ao
juntá-los: estados que esperam a mesma continuação, como os dois
algarismos finais, passaram a ser um só.

O trabalho também deixou claro o que o formalismo garante e o que
não garante. Ele garante resposta exata e previsível para o
formato, e permite provar que está correta. Ele não garante que a
placa exista, e as decisões sobre hífen, letras maiúsculas e
quinta posição do Mercosul precisaram ser tomadas e documentadas,
pois a norma sozinha não as resolve.

\subsection{Considerações finais}

Para quem vai usar o validador, o resultado prático é uma
verificação rápida e confiável do formato da placa, que reduz
erros de digitação antes de qualquer consulta a cadastros. Para a
equipe, o trabalho mostrou como a teoria de autômatos se aplica a
um problema do dia a dia, do modelo matemático ao código, e como
cada etapa pode ser verificada. As limitações e as possibilidades
de evolução estão na Seção~\ref{sec:limitacoes}.